% ============================================================
\chapter{Stationnarit\'e et Autocovariance}
\label{ch:stationnarite}
% ============================================================

\section{Exemple motivant}
La temp\'erature journali\`ere \`a Paris sur 30 ans pr\'esente une saisonnalit\'e marqu\'ee mais, apr\`es d\'esaisonnalisation, les r\'esidus semblent fluctuer autour d'une moyenne constante avec une structure de d\'ependance invariante. C'est la \textbf{stationnarit\'e}.

\section{Stationnarit\'e stricte et faible}

\begin{definition}[Stationnarit\'e stricte]\index{stationnarite@stationnarit\'e!stricte}
$(X_t)$ est \textbf{strictement stationnaire} si pour tout $k\geq 1$, tout $(t_1,\ldots,t_k)$ et tout $h\in\Z$ :
\[
(X_{t_1},\ldots,X_{t_k}) \stackrel{d}{=} (X_{t_1+h},\ldots,X_{t_k+h}).
\]
\end{definition}

\begin{definition}[Stationnarit\'e faible (au second ordre)]\index{stationnarite@stationnarit\'e!faible}
$(X_t)$ est \textbf{faiblement stationnaire} (ou stationnaire au sens large) si :
\begin{enumerate}
\item $\E[X_t^2]<\infty$ pour tout $t$,
\item $\mu = \E[X_t]$ ne d\'epend pas de $t$,
\item $\gamma(t,t+h) = \Cov(X_t,X_{t+h})$ ne d\'epend que de $h$.
\end{enumerate}
\end{definition}

\begin{remarque}
Stationnarit\'e stricte + moments d'ordre 2 finis $\Rightarrow$ stationnarit\'e faible. La r\'eciproque est fausse en g\'en\'eral, mais vraie pour les processus gaussiens.
\end{remarque}

\section{Fonction d'autocovariance}

\begin{definition}[Autocovariance et autocorr\'elation]\index{autocovariance}\index{autocorrelation@autocorr\'elation}
Pour un processus faiblement stationnaire :
\begin{align}
\gamma(h) &= \Cov(X_t,X_{t+h}) = \E[(X_t-\mu)(X_{t+h}-\mu)], \\
\rho(h) &= \frac{\gamma(h)}{\gamma(0)} = \Corr(X_t,X_{t+h}).
\end{align}
La fonction $\rho(\cdot)$ est la \textbf{fonction d'autocorr\'elation} (ACF)\index{ACF}.
\end{definition}

\begin{theoreme}[Propri\'et\'es de $\gamma$]\label{thm:gamma-prop}
\begin{enumerate}
\item $\gamma(0)\geq 0$ (c'est la variance).
\item $\gamma(h)=\gamma(-h)$ (sym\'etrie).
\item $|\gamma(h)|\leq\gamma(0)$ (par Cauchy--Schwarz).
\item $\gamma$ est \textbf{semi-d\'efinie positive} : pour tout $n$, tout $(a_1,\ldots,a_n)\in\R^n$ et tout $(t_1,\ldots,t_n)$ :
\[
\sum_{i,j}a_i a_j\gamma(t_i-t_j)\geq 0.
\]
\end{enumerate}
\end{theoreme}

\begin{proof}
(4) : $\sum_{i,j}a_ia_j\gamma(t_i-t_j)=\Var\!\left(\sum_i a_i X_{t_i}\right)\geq 0$.
\end{proof}

\begin{theoreme}[Herglotz]\index{Herglotz}
$\gamma:\Z\to\R$ est une fonction d'autocovariance si et seulement si elle est semi-d\'efinie positive. De plus, il existe une mesure positive finie $F$ sur $[-\pi,\pi]$ (mesure spectrale) telle que :
\[
\gamma(h) = \int_{-\pi}^{\pi} e^{i\omega h}\,dF(\omega).
\]
\end{theoreme}

\section{Estimation de l'autocovariance}

\begin{definition}[Autocovariance empirique]\index{autocovariance!empirique}
\[
\hat{\gamma}(h) = \frac{1}{n}\sum_{t=1}^{n-|h|}(X_t-\bar{X})(X_{t+|h|}-\bar{X}), \quad \hat{\rho}(h)=\frac{\hat{\gamma}(h)}{\hat{\gamma}(0)}.
\]
On divise par $n$ (et non $n-|h|$) pour garantir que la matrice $(\hat{\gamma}(i-j))$ est semi-d\'efinie positive.
\end{definition}

\begin{theoreme}[Bartlett]\index{Bartlett}
Pour un bruit blanc de taille $n$ : $\hat{\rho}(h)\stackrel{\text{approx}}{\sim}\mathcal{N}(0,1/n)$ pour $h\geq 1$. Les \textbf{bandes de Bartlett} $\pm 1{,}96/\sqrt{n}$ servent \`a tester si $\rho(h)=0$.
\end{theoreme}

\section{Autocorr\'elation partielle}\index{PACF}

\begin{definition}[PACF]
L'\textbf{autocorr\'elation partielle} au d\'ecalage $h$ est :
\[
\alpha(h) = \Corr(X_t-\hat{X}_t,\;X_{t+h}-\hat{X}_{t+h}),
\]
o\`u $\hat{X}_t$ et $\hat{X}_{t+h}$ sont les projections lin\'eaires sur $(X_{t+1},\ldots,X_{t+h-1})$.

\'Equivalemment, $\alpha(h)$ est le dernier coefficient $\phi_{hh}$ dans la r\'egression lin\'eaire de $X_t$ sur $(X_{t-1},\ldots,X_{t-h})$.
\end{definition}

\begin{theoreme}
Pour un AR$(p)$ : $\alpha(h)=0$ pour $h>p$. C'est l'outil principal pour identifier l'ordre $p$.
\end{theoreme}

\section{Tests de stationnarit\'e}\index{test!Dickey--Fuller}\index{test!KPSS}

\begin{definition}[Test de Dickey--Fuller augment\'e (ADF)]
Pour tester la pr\'esence d'une racine unitaire dans $X_t = \phi X_{t-1}+\eps_t$ :
\[
H_0 : \phi=1 \quad\text{(non stationnaire)} \quad \text{vs} \quad H_1 : |\phi|<1 \quad\text{(stationnaire)}.
\]
On estime $\nabla X_t = (\phi-1)X_{t-1}+\eps_t$ et on teste si $\phi-1<0$.
\end{definition}

\begin{attention}
La distribution de la statistique ADF n'est \textbf{pas} gaussienne sous $H_0$. On utilise les tables de Dickey--Fuller (ou les valeurs critiques calcul\'ees par simulation).
\end{attention}

\section{Impl\'ementation}

\begin{codeblock}[title=\textbf{Python}]
\begin{minted}{python}
import numpy as np
import matplotlib.pyplot as plt
from statsmodels.tsa.stattools import acf, pacf, adfuller
from statsmodels.graphics.tsaplots import plot_acf, plot_pacf

# Simuler un AR(1)
np.random.seed(0)
n = 500
phi = 0.7
eps = np.random.normal(0, 1, n)
x = np.zeros(n)
for t in range(1, n):
    x[t] = phi * x[t-1] + eps[t]

# ACF et PACF
fig, axes = plt.subplots(1, 2, figsize=(12, 4))
plot_acf(x, lags=30, ax=axes[0], title='ACF')
plot_pacf(x, lags=30, ax=axes[1], title='PACF')
plt.tight_layout()
plt.savefig('ch02_acf_pacf.pdf')
plt.show()

# Test ADF
result = adfuller(x, regression='c')
print(f"ADF statistic: {result[0]:.4f}")
print(f"p-value: {result[1]:.4f}")
print(f"Stationnaire: {result[1] < 0.05}")

# Marche aléatoire (non stationnaire)
rw = np.cumsum(eps)
result_rw = adfuller(rw, regression='c')
print(f"\nMarche aléatoire - ADF: {result_rw[0]:.4f}, p={result_rw[1]:.4f}")
\end{minted}
\end{codeblock}

\section{Exercices}

\begin{exercice}[$\star$]
Calculer la fonction d'autocovariance d'un bruit blanc $\eps_t\sim\mathrm{BB}(0,\sigma^2)$.
\end{exercice}

\begin{exercice}[$\star$]
Calculer l'ACF de $X_t = \eps_t + \theta\eps_{t-1}$ (MA(1)). Montrer que $\rho(h)=0$ pour $|h|\geq 2$.
\end{exercice}

\begin{exercice}[$\star\star$]
Montrer que si $(X_t)$ est strictement stationnaire avec $\E[X_t^2]<\infty$, alors $(X_t)$ est faiblement stationnaire.
\end{exercice}

\begin{exercice}[$\star\star$]
D\'emontrer que l'autocovariance empirique $\hat{\gamma}(h)$ (divis\'ee par $n$) donne une matrice de Toeplitz semi-d\'efinie positive.
\end{exercice}

\begin{exercice}[$\star\star\star$ -- Projet]
Appliquer le test ADF et le test KPSS \`a 10 s\'eries \'economiques (PIB, ch\^omage, inflation, etc.). Comparer les conclusions et discuter les cas de d\'esaccord.
\end{exercice}

\begin{formules}
\begin{align*}
\gamma(h) &= \Cov(X_t,X_{t+h}), \quad \rho(h)=\gamma(h)/\gamma(0) \\
\hat{\rho}(h) &\stackrel{\text{approx}}{\sim} \mathcal{N}(0,1/n) \quad\text{sous }H_0:\rho(h)=0 \\
\text{ADF} &: \nabla X_t = (\phi-1)X_{t-1}+\eps_t, \quad H_0:\phi=1
\end{align*}
\end{formules}
